% GENERATED FILE. Edit canonical lessons, then run npm run generate:books.
% canonical-source-hash: fnv1a64:8363f38aff54a726
% canonical-lessons: GE-C466-eifersucht, GE-C466-neid, GE-C466-neidisch, GE-C466-beneiden, GE-C466-goennen

\chapter{Neid und Eifersucht}
\label{ch:ge-neid-und-eifersucht}
\hlchaptermodality{466}

\begin{chapteropening}
\textbf{By the end of this chapter:} \emph{I can talk about envy and jealousy.}

Complete the last lesson of chapter 466: I can talk about envy and jealousy.
\end{chapteropening}

\section[die Eifersucht]{die Eifersucht — jealousy}

\label{lesson:GE-C466-eifersucht}



\begin{quote}
Before the new one: say the German for at least, then the German for during.
\end{quote}

\subsection*{You'll want to know: die Eifersucht}
\textbf{die Eifersucht} — \textquotedblleft{}jealousy\textquotedblright{}.

\begin{quote}
Seine Eifersucht macht seine Freundin traurig.
\end{quote}

\begin{quote}
\emph{His jealousy makes his girlfriend sad.}
\end{quote}

No plural in everyday use. The adjective is \emph{eifersüchtig} — jealous.

\begin{cousinweb}
\emph{Eifersucht} is built on \emph{Eifer} \textquotedblleft{}zeal, eagerness\textquotedblright{} and \emph{Sucht} \textquotedblleft{}sickness, craving\textquotedblright{}; \emph{Sucht} is inherited from Old High German \emph{suht}, from Proto-Germanic \emph{suhtiz} (EtymDB).
\end{cousinweb}

\begin{grammarlens}[title={What you've built}]
One more word for telling, explaining and describing at B1.
\end{grammarlens}

\subsection*{Your turn}
\begin{itemize}
\raggedright
  \item \emph{Say it:} \emph{die Eifersucht}
  \item \emph{Say it:} \emph{die Eifersucht}, once more
  \item \emph{Say it:} a sentence of your own with \emph{die Eifersucht}
  \item \emph{Recall:} say the German for at least, then the German for during, then say \emph{die Eifersucht} again
\end{itemize}

\begin{tcolorbox}[breakable,colback=teal!4,colframe=teal!35!black,title={Before you move on}]
What does die Eifersucht mean? (\textquotedblleft{}Jealousy\textquotedblright{}.) Say it once more.
\end{tcolorbox}
\section[der Neid]{der Neid — envy}

\label{lesson:GE-C466-neid}



\begin{quote}
Before the new one: say the German for during, then the German for jealousy.
\end{quote}

\subsection*{You'll want to know: der Neid}
\textbf{der Neid} — \textquotedblleft{}envy\textquotedblright{}.

\begin{quote}
Er spricht ohne Neid über den Erfolg seines Bruders.
\end{quote}

\begin{quote}
\emph{He talks about his brother's success without envy.}
\end{quote}

No plural. \emph{Vor Neid platzen} — to burst with envy.

\begin{cousinweb}
\emph{Neid} is a native Germanic word, usually traced to Old High German \emph{nīd} \textquotedblleft{}hatred, envy\textquotedblright{}; EtymDB lists the German word without a source.
\end{cousinweb}

\begin{grammarlens}[title={What you've built}]
One more word for telling, explaining and describing at B1.
\end{grammarlens}

\subsection*{Your turn}
\begin{itemize}
\raggedright
  \item \emph{Say it:} \emph{der Neid}
  \item \emph{Say it:} \emph{der Neid}, once more
  \item \emph{Say it:} a sentence of your own with \emph{der Neid}
  \item \emph{Recall:} say the German for during, then the German for jealousy, then say \emph{der Neid} again
\end{itemize}

\begin{tcolorbox}[breakable,colback=teal!4,colframe=teal!35!black,title={Before you move on}]
What does der Neid mean? (\textquotedblleft{}Envy\textquotedblright{}.) Say it once more.
\end{tcolorbox}
\section[neidisch]{neidisch — envious}

\label{lesson:GE-C466-neidisch}



\begin{quote}
Before the new one: say the German for jealousy, then the German for envy.
\end{quote}

\subsection*{You'll want to know: neidisch}
\textbf{neidisch} — \textquotedblleft{}envious\textquotedblright{}.

\begin{quote}
Meine Schwester ist neidisch auf mein neues Fahrrad.
\end{quote}

\begin{quote}
\emph{My sister is envious of my new bike.}
\end{quote}

\emph{Neidisch auf} + accusative — envious of.

\begin{cousinweb}
\emph{neidisch} is built on \emph{Neid} \textquotedblleft{}envy\textquotedblright{} with the adjective ending \emph{-isch}; \emph{Neid} is usually traced to Old High German \emph{nīd} \textquotedblleft{}hatred, envy\textquotedblright{} (EtymDB lists both German words).
\end{cousinweb}

\begin{grammarlens}[title={What you've built}]
One more word for telling, explaining and describing at B1.
\end{grammarlens}

\subsection*{Your turn}
\begin{itemize}
\raggedright
  \item \emph{Say it:} \emph{neidisch}
  \item \emph{Say it:} \emph{neidisch}, once more
  \item \emph{Say it:} a sentence of your own with \emph{neidisch}
  \item \emph{Recall:} say the German for jealousy, then the German for envy, then say \emph{neidisch} again
\end{itemize}

\begin{tcolorbox}[breakable,colback=teal!4,colframe=teal!35!black,title={Before you move on}]
What does neidisch mean? (\textquotedblleft{}Envious\textquotedblright{}.) Say it once more.
\end{tcolorbox}
\section[beneiden]{beneiden — to envy}

\label{lesson:GE-C466-beneiden}



\begin{quote}
Before the new one: say the German for envy, then the German for envious.
\end{quote}

\subsection*{You'll want to know: beneiden}
\textbf{beneiden} — \textquotedblleft{}to envy\textquotedblright{}.

\begin{quote}
Ich beneide dich um deinen ruhigen Job.
\end{quote}

\begin{quote}
\emph{I envy you your quiet job.}
\end{quote}

\emph{Jemanden um etwas beneiden} — to envy someone something. Inseparable: \emph{hat beneidet}.

\begin{cousinweb}
\emph{beneiden} is built on the prefix \emph{be-} and \emph{Neid} \textquotedblleft{}envy\textquotedblright{}, a Germanic word usually traced to Old High German \emph{nīd} (EtymDB lists the verb without a source).
\end{cousinweb}

\begin{grammarlens}[title={What you've built}]
One more word for telling, explaining and describing at B1.
\end{grammarlens}

\subsection*{Your turn}
\begin{itemize}
\raggedright
  \item \emph{Say it:} \emph{beneiden}
  \item \emph{Say it:} \emph{beneiden}, once more
  \item \emph{Say it:} a sentence of your own with \emph{beneiden}
  \item \emph{Recall:} say the German for envy, then the German for envious, then say \emph{beneiden} again
\end{itemize}

\begin{tcolorbox}[breakable,colback=teal!4,colframe=teal!35!black,title={Before you move on}]
What does beneiden mean? (\textquotedblleft{}To envy\textquotedblright{}.) Say it once more.
\end{tcolorbox}
\section[gönnen]{gönnen — to not begrudge; to treat oneself (sich etwas gönnen)}

\label{lesson:GE-C466-goennen}



\begin{quote}
Before the new one: say the German for envious, then the German for to envy.
\end{quote}

\subsection*{You'll want to know: gönnen}
\textbf{gönnen} — \textquotedblleft{}to not begrudge; to treat oneself (sich etwas gönnen)\textquotedblright{}.

\begin{quote}
Ich gönne dir den Erfolg von Herzen.
\end{quote}

\begin{quote}
\emph{I really don't begrudge you your success.}
\end{quote}

\emph{Sich etwas gönnen} — to treat yourself: \emph{Heute gönne ich mir ein Stück Kuchen.}

\begin{cousinweb}
\emph{gönnen} is inherited from Old High German \emph{giunnan}, built with \emph{gi-} on \emph{unnan} \textquotedblleft{}to grant\textquotedblright{}, from Proto-Germanic \emph{unnaną} \textquotedblleft{}to grant, to bestow\textquotedblright{} (EtymDB).
\end{cousinweb}

\begin{grammarlens}[title={What you've built}]
One more word for telling, explaining and describing at B1.
\end{grammarlens}

\subsection*{Your turn}
\begin{itemize}
\raggedright
  \item \emph{Say it:} \emph{gönnen}
  \item \emph{Say it:} \emph{gönnen}, once more
  \item \emph{Say it:} a sentence of your own with \emph{gönnen}
  \item \emph{Recall:} say the German for envious, then the German for to envy, then say \emph{gönnen} again
\end{itemize}

\begin{tcolorbox}[breakable,colback=teal!4,colframe=teal!35!black,title={Before you move on}]
What does gönnen mean? (\textquotedblleft{}To not begrudge; to treat oneself (sich etwas gönnen)\textquotedblright{}.) Say it once more.
\end{tcolorbox}
